\begin{abstract}
We study cross-lingual named entity recognition (NER) transfer to Danish using two multilingual transformer encoders — mBERT and XLM-R — pitting English against Norwegian as source languages across three random seeds. On the held-out Danish test set, Norwegian comes out ahead in the zero-shot setting: +1.86 F1 for XLM-R ($p=0.10$) and +3.36 F1 for mBERT ($p=0.046$, significant at $\alpha=0.05$). One run — seed 42 — looked especially favorable to Norwegian XLM-R on the development set, then revealed itself as an outlier once averaged across seeds, which is precisely why cross-lingual source-language comparisons should not rest on a single initialization. The results track what typological structure would lead you to expect: Norwegian is roughly 5.7$\times$ closer to Danish than English in WALS syntactic-feature space (cosine distance 0.021 vs.\ 0.119). The picture changes once Danish fine-tuning enters: an English mBERT checkpoint achieves the best transfer-plus-fine-tune dev score (89.4 strict F1), suggesting that a large source corpus can offset typological distance once target-language supervision is available. We additionally report unlabeled and loose span F1, per-entity-type breakdowns, a Danish data-efficiency curve, and a diagnosis of one misaligned prediction file.
\end{abstract}
